These entries use the Greeks to hedge options: delta and gamma hedging, hedging when implied and actual volatility differ, and hedging with transaction costs. Numbers are checked in \texttt{notes/artefacts/option-valuation/option-valuation-checks.py}; rates are continuously compounded unless stated.

\begin{qfmethod}[The Greeks in use: interpreting them, executing and rebalancing a delta hedge, gamma risk and gamma-neutral hedging, and implied volatility in trading]\label{meth:greeks-delta-gamma-hedging}
\qffield{Statement} \emph{Interpretation.}
\begin{itemize}
\item Delta is the change in option value per unit change in the underlying.
\item Gamma is the change in delta.
\item Theta is the change with the passage of time.
\item Vega is the change per unit of volatility.
\item Rho is the change per unit of interest rate.
\end{itemize}
Formulas are in \hyperref[mod:black-scholes-merton]{the BSM model}. A position with zero delta is delta neutral, but only for a short time, so it must be rebalanced: \emph{dynamic hedging} \citep[chap.~19, sec.~19.4, pp.~421--427]{s7-hulloptions}.

\emph{Executing a delta hedge.} A bank that sells calls buys $N(d_1)$ shares per call and adjusts the holding as delta changes. If the hedge worked perfectly, the discounted cost of writing and hedging the option would equal its BSM price on every path; weekly rebalancing leaves a dispersion that falls as rebalancing becomes more frequent (Tables~19.2--19.4, pp.~424--426).

\emph{Gamma risk.} For a delta-neutral portfolio $\Delta\Pi\approx\Theta\Delta t+\tfrac12\Gamma\Delta S^2$ (eq.~(19.3), p.~430). With negative gamma the portfolio gains if the price does not move and loses on large moves of either sign. The underlying has zero gamma, so a portfolio is made gamma neutral with $-\Gamma/\Gamma_T$ traded options and then delta neutral again with the underlying (pp.~430--431). Vega is hedged the same way, with options (sec.~19.8).

\emph{Second-order volatility Greeks.} Vanna, $\partial^2V/\partial S\partial\sigma$, is the change in delta per unit of volatility; volga, $\partial^2V/\partial\sigma^2$, is the change in vega per unit of volatility. Wilmott records them under the names ``kabanga'' and ``vomma'' and warns that vega itself is meaningful only for payoffs whose gamma has one sign \citep[chap.~7, pp.~127--128 and n.~3]{wide-wilmott-pwoqf}. For a European call, vanna is $-e^{-qT}N'(d_1)d_2/\sigma$ and volga is $\text{vega}\cdot d_1d_2/\sigma$.

\emph{Implied volatility in trading.} The implied volatility is the volatility that makes the model price equal the market price. Traders quote options in implied volatility, and the VIX indexes the 30-day implied volatility of S\&P~500 options (chap.~15, sec.~15.11, pp.~358--359). A delta-hedged long option earns, roughly, $\tfrac12\Gamma S^2(\sigma_{\text{realised}}^2-\sigma_{\text{implied}}^2)\Delta t$ each period. So buying options is a view that realised volatility will exceed implied: ``If implied volatility is low, such that options generally appear underpriced, look for spreads with a positive vega'' \citep[chap.~11, p.~199]{s7-natenberg}.

\qffield{Practice} \emph{Reproduced} (Hull, Example~19.1 and sec.~19.6).
\begin{itemize}
\item A call with $S=49$, $K=50$, 5 percent, 20 percent volatility and 20 weeks has a delta of 0.522; 100{,}000 calls are worth about \$240{,}000.
\item A delta-neutral book with gamma $-3{,}000$ and a traded call with delta 0.62 and gamma 1.50 needs 2{,}000 calls bought and 1{,}240 shares sold.
\item A book with gamma $-10{,}000$ loses about 20{,}000 on a move of 2.
\end{itemize}

\emph{Checked} (original, the Example~19.1 call).
\begin{itemize}
\item \emph{Rebalancing frequency.} Simulating 20{,}000 paths with a 13 percent drift, the ratio of the standard deviation of hedging cost to the BSM price is 0.415, 0.373, 0.270 and 0.192 for rebalancing every 5, 4, 2 and 1 weeks, within 0.03 of Hull's 0.42, 0.38, 0.28 and 0.21. The mean cost, 2.4035, is close to the 2.40 price.
\item \emph{Short intervals.} At half a week and a quarter of a week the simulation gives 0.136 and 0.098, falling by about $\sqrt2$ per halving as discrete-hedging theory predicts. Hull's Table~19.4 gives 0.16 and 0.13, ratios of 1.31 and 1.23; the cause is not stated, and the discrepancy is recorded.
\item \emph{Gamma--theta.} A short call hedged with delta shares over one day changes by $-0.1176$ for a move of $-2$ and $-0.1095$ for $+2$ on full revaluation, against $-0.1140$ from $\Theta\Delta t+\tfrac12\Gamma\Delta S^2$.
\item \emph{Implied against realised.} Buying the call at 20 percent implied and hedging daily at that volatility, while realised volatility is 25 percent, earns 0.62 on average per call (standard error 0.006).
\item \emph{Vega.} Vega is 12.1 per unit of volatility, so one point of implied volatility moves the call by about 0.12.
\end{itemize}

\qffield{Conditions} Constant volatility, no transaction costs and continuous trading make the hedge perfect in theory. In practice traders rebalance to delta daily and manage gamma and vega by watching them, because options to neutralise them are expensive, and they use scenario tables of profit and loss across prices and volatilities (secs.~19.10--19.11, pp.~437--438). Greeks computed at one implied volatility per option ignore smile dynamics (\hyperref[def:implied-volatility-surface]{volatility surface}).

\qffield{Pitfalls} Treating a delta-neutral book as riskless. Hedging gamma with the underlying. Reading a profit on a delta-hedged option as alpha when it is the realised-minus-implied variance spread, earned at the cost of gamma risk.

\qffield{See also} \hyperref[mod:black-scholes-merton]{BSM}; \hyperref[meth:option-spreads-volatility-strategies]{volatility strategies}; \hyperref[def:vix-variance-dividend-swaps]{the VIX}; \hyperref[res:delta-hedging-volatility-mismatch]{hedging with the wrong volatility} (which delta to use when implied and actual volatility differ).

\end{qfmethod}

\begin{qfresult}[Delta hedging when implied and actual volatility differ: the locked-in profit of the actual-volatility hedge against the path-dependent profit of the implied-volatility hedge]\label{res:delta-hedging-volatility-mismatch}
\qffield{Statement} Suppose an option is bought at the price implied by $\tilde\sigma$ while the asset's actual volatility is $\sigma$, and the position is delta hedged to expiry. The question is which volatility goes into the delta \citep[chap.~12, pp.~197--199]{wide-wilmott-pwoqf}.
\begin{itemize}
\item \emph{Hedge with actual volatility.} This replicates a short position in a correctly priced option, so the profit is guaranteed to be $V(S,t;\sigma)-V(S,t;\tilde\sigma)$ in present value by expiry. The mark-to-market profit over each step is $\tfrac12(\sigma^2-\tilde\sigma^2)S^2\Gamma^i\,dt+(\Delta^i-\Delta^a)((\mu-r+D)S\,dt+\sigma S\,dX)$: it contains a $dX$ term and the real drift $\mu$, so ``you could lose before you gain'' (pp.~200--202).
\item \emph{Hedge with implied volatility.} The profit each step is $\tfrac12(\sigma^2-\tilde\sigma^2)S^2\Gamma^i\,dt$, eq.~(12.1), with no $dX$ term. Its present value over the life is $\tfrac12(\sigma^2-\tilde\sigma^2)\int_{t_0}^Te^{-r(t-t_0)}S^2\Gamma^i\,dt$, which is ``always positive, but highly path dependent''; the positivity needs $\Gamma^i\ge0$, as for a vanilla call or put. It depends on $\mu$, is largest if the asset ends near the strike, and needs no estimate of actual volatility beyond its sign relative to implied (pp.~202--203). For a payoff whose gamma changes sign the integrand changes sign too, so the total can be negative even with $\sigma>\tilde\sigma$.
\end{itemize}
Hedging with any volatility $\sigma_h$ gives $V(S,t;\sigma_h)-V(S,t;\tilde\sigma)+\tfrac12(\sigma^2-\sigma_h^2)\int e^{-r(t-t_0)}S^2\Gamma^h\,dt$, eq.~(12.2), after Carr and Henrard (p.~203).

\qffield{Notation} Superscripts $a$ and $i$ for values and Greeks at actual and implied volatility; $D$ dividend yield; profits in present value at $t_0$. Interpretation: ``zero diffusion'' describes the continuous-time increment of the implied-volatility hedge. With rebalancing at intervals $\delta t$ both hedges carry a hedging error, which Wilmott writes as $\tfrac12(\phi^2-1)\sigma^2S^2\Gamma\,\delta t$ per step: mean zero, chi-squared and skewed \citep[chap.~47, pp.~764--770]{wide-wilmott-pwoqf}. Hull's simulations show the dispersion of a discretely rebalanced hedge falling steadily with the rebalancing interval \citep[chap.~19, sec.~19.4, pp.~425--426]{s7-hulloptions}.

\qffield{Conditions} Geometric Brownian motion with constant, known actual volatility and constant implied volatility; the option is marked at implied volatility throughout; continuous rebalancing, no costs, borrowing and lending at $r$. The guaranteed profit needs the actual volatility to be known and constant, which it never is; the implied-volatility hedge has a positive total if $\sigma>\tilde\sigma$ throughout and $\Gamma^i\ge0$ along the path, as for a vanilla call or put (a convex payoff). Neither profit is positive merely because actual volatility exceeds implied: with negative gamma the integrand is negative, and where vega is negative $V^a-V^i<0$. With a smile, stochastic volatility or jumps, neither result holds exactly. The examples are synthetic.

\qffield{Practice} Example, synthetic: a one-year at-the-money call, $S_0=K=100$, $r=5$ percent, $q=0$, real drift 10 percent, actual volatility 30 percent, bought at implied 20 percent. Then $V^a-V^i=14.2313-10.4506=3.7807$. On 20{,}000 common paths:
\begin{itemize}
\item \emph{Actual-volatility delta.} Mean profit 3.7770 (standard error 0.0044) with 252 rebalancings and 3.7805 (0.0014) with 2{,}520. The standard deviation falls from 0.627 to 0.200, a ratio of 3.14 against $\sqrt{10}=3.16$: the dispersion is discrete-hedging error.
\item \emph{Implied-volatility delta.} The root-mean-square gap between each path's profit and its pathwise integral falls from 0.797 to 0.256. The standard deviation across paths stays at 1.75, because the total itself depends on the path. The mean profit here is 3.746 (standard error 0.012); with another drift it would differ.
\item \emph{The gamma condition, in closed form.} A unit cash-or-nothing call (\hyperref[meth:digital-option-valuation]{digital options}) with $S_0=120$, $K=100$, $r=5$ percent, $q=0$ and $T=1$ is worth 0.8141 at 20 percent volatility and 0.6980 at 30 percent, so the actual-volatility hedge locks in $-0.1160$. Its gamma at implied volatility is $+0.00071$ at $S=80$ and $-0.00047$ at $S=120$; it is negative wherever $d_1>0$, so a path that stays above the strike makes the implied-volatility hedge lose.
\end{itemize}

\qffield{Pitfalls} Expecting the implied-volatility hedge to earn $V^a-V^i$, or the actual-volatility hedge to make money smoothly. Calling the implied-volatility hedge riskless because each increment has no $dW$ term: the total is random, and finite rebalancing adds noise. Forgetting that the guaranteed profit depends on a volatility forecast that has to be right for the whole life. Expecting a positive variance spread to pay on a digital, barrier or spread position whose gamma changes sign.

\qffield{See also} \hyperref[meth:greeks-delta-gamma-hedging]{the Greeks in use}; \hyperref[mod:black-scholes-merton]{BSM}; \hyperref[def:vix-variance-dividend-swaps]{variance swaps} (trading the variance spread directly); \hyperref[meth:active-currency-strategies]{volatility trading in currencies}; \hyperref[meth:transaction-cost-hedging]{hedging with transaction costs}.

\end{qfresult}

\begin{qfmethod}[Hedging options with proportional transaction costs: Leland's adjusted volatility for long and short positions, the nonlinear portfolio equation, and fixed-interval against band rebalancing]\label{meth:transaction-cost-hedging}
\qffield{Statement} With a bid-offer spread on the underlying, continuous rebalancing would cost an unbounded amount, so a hedger rebalances at intervals and pays for each trade. The option then has no single value: costs are always a drain, so a hedger values a long position below the Black--Scholes value and a short position above it, and the amounts depend on the hedger's own costs. The modelling was initiated by Leland (1985). In Leland's model the portfolio is rebalanced to the delta every $\delta t$; buying or selling $\nu$ shares at $S$ costs $\kappa|\nu|S$; and the hedged portfolio's expected return is set equal to the risk-free rate. Long calls and puts are then valued with $\check\sigma^2=\sigma^2-2\kappa\sigma\sqrt{2/(\pi\delta t)}$ and short ones with $\hat\sigma^2=\sigma^2+2\kappa\sigma\sqrt{2/(\pi\delta t)}$, eqs.~(48.5)--(48.6) \citep[chap.~48, secs.~48.1--48.4, pp.~783--789]{wide-wilmott-pwoqf}.

The derivation explains the signs. Over one step the shares traded are, to leading order, $\nu\approx\Gamma\,\delta S=\Gamma\sigma S\phi\sqrt{\delta t}$, so the expected cost is $\sqrt{2/\pi}\,\kappa\sigma S^2|\Gamma|\sqrt{\delta t}$, eq.~(48.2). Subtracting it from the portfolio's expected return gives the Hoggard--Whalley--Wilmott equation $V_t+\tfrac12\sigma^2S^2V_{SS}-\kappa\sigma S^2\sqrt{2/(\pi\delta t)}\,|V_{SS}|+rSV_S-rV=0$, eq.~(48.4). For a single long vanilla option $\Gamma>0$, the modulus drops and the equation is Black--Scholes at $\check\sigma$. For a short one every sign changes except the cost term's, giving $\hat\sigma$. A hedger selling after a rise sells at the bid, so the effective price moves less than the market: the long position sees a lower volatility (pp.~786--788; appendix, pp.~809--811).

The equation is nonlinear, so a portfolio's value is not the sum of its parts. Long and short the same call, hedged separately, lose on both hedges, while hedged together they never trade. For a portfolio whose gamma changes sign the modulus cannot be dropped and the equation must be solved numerically. The marginal effect of adding a contract depends on the sign of the existing portfolio's gamma: adding gamma of the opposite sign reduces costs (secs.~48.5--48.6 and 48.10.1, pp.~790--793 and 797).

\qffield{Notation} $\kappa$ one-way proportional cost; Interpretation: if trades fill at mid plus or minus half the spread, $\kappa$ is half the proportional bid-offer spread. $\delta t$ the rebalancing interval; Wilmott's non-dimensional cost $\mathcal K=\kappa/(\sigma\sqrt{\delta t})$, eq.~(48.7), written $K$ there. Large $\mathcal K$ means rehedging too often for the costs, small $\mathcal K$ too rarely (pp.~788--789). For small costs Wilmott's approximation $2\kappa SN'(d_1)\sqrt{T-t}/\sqrt{2\pi\delta t}$, the volatility change times vega, is each side's adjustment from the frictionless value: the short value lies about that much above Black--Scholes and the long value about that much below (p.~788). Interpretation: the full short-minus-long difference is therefore about twice that amount to first order. Interpretation: by It\^o's formula along paths of volatility $\sigma$, exactly $\hat V-V=E^Q\int_0^Te^{-rt}\kappa\sigma\sqrt{2/(\pi\delta t)}\,S_t^2\hat\Gamma_t\,dt$, where $\hat V$ and $\hat\Gamma$ are at $\hat\sigma$. The Leland premium is therefore the expected present value of the leading-order cost of hedging with the $\hat\sigma$ delta; the long discount $V-\check V$ is the same with $\check\Gamma$.

\qffield{Conditions} Leland's assumptions are rebalancing at a fixed, finite, small interval, costs proportional to the value traded, normal returns and valuation by the hedged portfolio's expected return. Choosing $\Delta=\partial V/\partial S$ minimises the leading-order variance but leaves a residual risk. The adjusted volatility applies only to positions whose gamma has one sign. The long formula needs $\check\sigma^2>0$, which by (48.5) means $\kappa<\sigma\sqrt{\pi\delta t/8}$, or $\mathcal K<\sqrt{\pi/8}\approx0.63$. Beyond that the equation is forward parabolic with final data and so ill-posed. Wilmott's p.~789 prints the threshold as $\kappa>\sigma/\sqrt{8\pi\delta t}$, a form that is dimensionally inconsistent with (48.5); the condition above follows from (48.5) directly. The formulas omit the trade that opens the hedge and any closing trade, and they are leading-order in $\sqrt{\delta t}$. For non-normal returns what matters is $E|\psi|$ of the standardised return; Wilmott's factors for selected assets are 0.50 to 0.96 against 1 for the normal (Table~48.1, p.~800). With fixed costs per trade the model can give negative values unless hedging stops below a boundary (sec.~48.10.2, p.~798). Wilmott's empirical comparisons use daily data for 1985--1996.

\qffield{Practice} \emph{Fixed interval against a band.} Costs fall and hedging risk rises as the interval lengthens: expected costs scale like $1/\sqrt{\delta t}$ and the hedging error's standard deviation like $\sqrt{\delta t}$, so a confidence-level price has an optimal interval (sec.~48.13, pp.~807--808). A delta-tolerance rule watches the price continuously and rebalances to the delta only when $\sigma S|D-\partial V/\partial S|$ exceeds a bandwidth $H_0$, eq.~(48.9), giving a different nonlinear equation, (48.10). Utility-based models (Hodges--Neuberger; Davis--Panas--Zariphopoulou) find an optimal no-trade band around the delta and trade only to its edge. Whalley and Wilmott's small-cost asymptotics make the band symmetric, with a width proportional to the cube root of the cost coefficient (secs.~48.8--48.9, pp.~793--797).

In Jefferies' tests of an at-the-money one-year call (volatility 20 percent, rate 5 percent, $S=100$, each cost coefficient 0.01), the lowest 95th-percentile prices on simulated paths were 14.90 for Leland hedging every four days, 15.03 for the utility band and 15.15 for delta tolerance. On real daily data for nearly 40 stocks, delta tolerance gave the lowest cost on about 70 percent of paths, the utility band on 23 percent and Leland on 7. Wilmott concludes that it is ``not at all clear which is the best strategy'' (sec.~48.12, pp.~800--806).

Examples, synthetic: an at-the-money call with $S_0=K=100$, $\sigma=20$ percent, $r=5$ percent, $q=0$ and $T=1$. Risk-neutral paths, 100{,}000, physical settlement, opening trade excluded; standard errors of simulated costs 0.0007 to 0.005.
\begin{itemize}
\item \emph{Leland premium against simulated cost.} At $\mathcal K=0.1$, $\hat\sigma=21.54$ and $\check\sigma=18.33$ percent. The short premium $\hat V-V$ is 0.5777 and the long discount $V-\check V$ 0.6233. Keeping $\mathcal K$ fixed with $\kappa=0.555$, 0.277 and 0.139 percent, the simulated present value of the $\hat\sigma$ hedge's costs over 13, 52 and 208 rebalances is 0.6177, 0.5902 and 0.5816. These exceed the premium by 6.9, 2.2 and 0.7 percent, a residual that shrinks about threefold per fourfold refinement. For the long position the residuals are 5.1, 1.6 and 0.5 percent; at $\mathcal K=0.4$, 9.5, 3.0 and 0.9 (short) and 1.9, 0.6 and 0.2 (long).
\item \emph{What the formula leaves out.} At the Leland price the writer's expected loss equals that residual. Opening the weekly hedge costs another 0.18, more than the residual. The settlement trade costs 0.017, against 0.18 for liquidating the hedge in a cash-settled contract.
\item \emph{Fixed cost, varying frequency.} At $\kappa=0.2$ percent with the Black--Scholes delta, 13, 52 and 208 rebalances cost 0.23, 0.44 and 0.87 and leave a profit-and-loss standard deviation of 1.77, 0.93 and 0.58. The premium that loses only 5 percent of the time is 13.7, 12.5 and 12.4, against a Black--Scholes value of 10.45.
\item \emph{Netting, a hedge-cost demonstration.} A bull spread, long the 95 call and short the 105 call, hedged weekly at $\kappa=0.2$ percent with Black--Scholes deltas: hedging the legs separately costs 0.854 in present value and trading only the net delta 0.247, a saving of 0.607 (standard error 0.0007). Netting never costs more on any step. Long and short the same call cost twice one hedge separately and nothing together. Valuing the legs separately with Leland's volatilities gives 4.49, against 5.33 without costs and 5.08 for the Black--Scholes spread less the netted cost. The last is an approximate value under a fixed delta rule, not a solution of (48.4).
\end{itemize}

\qffield{Pitfalls} Using the raised volatility for a long position, or the lowered one for a short. Applying the adjusted volatility to a spread or book whose gamma changes sign. Valuing a book as the sum of separately cost-adjusted options. Rehedging as often as possible, or using the long formula with $\mathcal K\ge\sqrt{\pi/8}$. Forgetting that the formula excludes the opening and closing trades.

\qffield{See also} \hyperref[res:delta-hedging-volatility-mismatch]{hedging with the wrong volatility} (no costs); \hyperref[meth:greeks-delta-gamma-hedging]{the Greeks in use}; \hyperref[mod:black-scholes-merton]{BSM}; \hyperref[res:option-market-maker-spreads]{option market-maker spreads}; \hyperref[def:spread-measures]{spread measures}.

\qffield{Further reading} \nolinkurl{vo-leland-1985} (the original model); \nolinkurl{vo-whalley-wilmott-1997} (the asymptotic optimal band).

\end{qfmethod}
